\chapter{Towards the continuum}
\label{chap:continuum}

The lattice arguments, written for a configuration space $(X,\mu)$ that may be continuous. States are real functions
in a form domain $D$, the potential-free energy (kinetic energy plus interaction) is a symmetric bilinear form $h_0$ on
$D$, and a one-body potential $w$ acts by multiplication with $Lw$; for $N$ particles in $\mathbb R^3$,
$Lw(x_1,\dots,x_N)=\sum_iw(x_i)$, and on the lattice $Lw(c)=\sum_iw_in_i(c)$. The density is the functional
$w\mapsto\int(Lw)\psi^2$. Analysis enters through three named hypotheses: the form domain separates functions
(fundamental lemma of the calculus of variations), the coupling is rigid, and the ground state satisfies unique
continuation ($\psi\neq0$ almost everywhere).

\begin{definition}[Setting]
  \label{def:continuum_setting}
  \lean{LatticeDFT.Continuum.Setting, LatticeDFT.Continuum.ip, LatticeDFT.Continuum.pf, LatticeDFT.Continuum.Q,
    LatticeDFT.Continuum.dens, LatticeDFT.Continuum.IsGS, LatticeDFT.Continuum.Sep, LatticeDFT.Continuum.RigidL}
  \leanok
  Form domain, energy form, coupling, ground states as minimisers of the Rayleigh quotient, separation and rigidity.
\end{definition}

\begin{lemma}[Euler--Lagrange equation]
  \label{lem:euler_lagrange}
  \lean{LatticeDFT.Continuum.integral_expand, LatticeDFT.Continuum.ip_comm, LatticeDFT.Continuum.pf_comm,
    LatticeDFT.Continuum.ip_expand, LatticeDFT.Continuum.pf_expand, LatticeDFT.Continuum.h0_expand,
    LatticeDFT.Continuum.euler_lagrange}
  \leanok
  \uses{def:continuum_setting}
  A ground state satisfies the weak Schr\"odinger equation $h_0(\psi,\varphi)+\langle w\psi,\varphi\rangle
  =E\langle\psi,\varphi\rangle$ for all $\varphi\in D$.
\end{lemma}
\begin{proof}
  \leanok
  The Rayleigh quotient along $\psi+s\varphi$ has a minimum at $s=0$, so its linear coefficient vanishes.
\end{proof}

\begin{theorem}[Hohenberg--Kohn, conditional on unique continuation]
  \label{thm:hk_continuum}
  \lean{LatticeDFT.Continuum.weak_hk, LatticeDFT.Continuum.strong_hk_of_ucp, LatticeDFT.Continuum.hk_of_ucp}
  \leanok
  \uses{lem:euler_lagrange}
  If ground states of $v$ and $v'$ have the same norm and density, they are ground states of both potentials; if one
  of them vanishes only on a null set, then $v-v'$ is constant.
\end{theorem}
\begin{proof}
  \leanok
  The usual variational argument gives the weak statement. Subtracting the two Euler--Lagrange equations,
  $\int(L(v-v')-c)\psi\varphi=0$ for every $\varphi\in D$; separation gives $(L(v-v')-c)\psi=0$ a.e., unique
  continuation gives $L(v-v')=c$ a.e., and rigidity gives a constant.
\end{proof}

\begin{theorem}[The lattice as an instance]
  \label{thm:hk_lattice_instance}
  \lean{LatticeDFT.Continuum.matForm, LatticeDFT.Continuum.latticeL, LatticeDFT.Continuum.latticeSetting,
    LatticeDFT.Continuum.ip_lattice, LatticeDFT.Continuum.Q_lattice, LatticeDFT.Continuum.isGS_lattice_iff,
    LatticeDFT.Continuum.sep_lattice, LatticeDFT.Continuum.rigidL_lattice, LatticeDFT.Continuum.ucp_lattice_iff,
    LatticeDFT.Continuum.lattice_hk_of_hstar}
  \leanok
  \uses{thm:hk_continuum}
  With counting measure on the configurations, separation is automatic, rigidity of the coupling is lattice rigidity,
  and unique continuation means no zero amplitude, which the Krylov condition (H*) guarantees. Hence Hohenberg--Kohn
  holds at every potential satisfying (H*).
\end{theorem}
\begin{proof}
  \leanok
  Test with indicator functions; (H*) gives full support of ground states.
\end{proof}

\section{One particle in a box: an unconditional theorem}

One particle on $[0,1]$ with hard walls, kinetic energy $\int\psi'^2$, form domain the $C^2$ functions vanishing at
$0$ and $1$, and continuous potentials on $[0,1]$. Every hypothesis of Theorem~\ref{thm:hk_continuum} is proved.

\begin{lemma}[The box setting]
  \label{lem:box_setting}
  \lean{LatticeDFT.Continuum.Box.μ1, LatticeDFT.Continuum.Box.Dom, LatticeDFT.Continuum.Box.mem_Dom,
    LatticeDFT.Continuum.Box.Pot, LatticeDFT.Continuum.Box.ext, LatticeDFT.Continuum.Box.continuous_ext,
    LatticeDFT.Continuum.Box.Lbox, LatticeDFT.Continuum.Box.integrable_of_continuous, LatticeDFT.Continuum.Box.kin,
    LatticeDFT.Continuum.Box.two_ge_one, LatticeDFT.Continuum.Box.diff_of_mem,
    LatticeDFT.Continuum.Box.cont_deriv_of_mem, LatticeDFT.Continuum.Box.setting,
    LatticeDFT.Continuum.Box.eq_zero_of_integral, LatticeDFT.Continuum.Box.sep_box, LatticeDFT.Continuum.Box.one,
    LatticeDFT.Continuum.Box.rigidL_box}
  \leanok
  \uses{def:continuum_setting}
  The box is a setting in which separation (fundamental lemma of the calculus of variations, with smooth test
  functions supported in $(0,1)$) and rigidity (a continuous potential that is a.e.\ constant is constant) hold.
\end{lemma}
\begin{proof}
  \leanok
  Mathlib's fundamental lemma on open sets; continuity and density of $(0,1)$ in $[0,1]$.
\end{proof}

\begin{lemma}[Ground states solve the Schr\"odinger equation]
  \label{lem:box_ode}
  \lean{LatticeDFT.Continuum.Box.energy, LatticeDFT.Continuum.Box.ode_of_isGS}
  \leanok
  \uses{lem:box_setting, lem:euler_lagrange}
  A ground state satisfies $\psi''=(v-E)\psi$ on $(0,1)$.
\end{lemma}
\begin{proof}
  \leanok
  Integrate the weak equation by parts (the boundary terms vanish) and apply the fundamental lemma.
\end{proof}

\begin{theorem}[Unique continuation in the box]
  \label{thm:box_ucp}
  \lean{LatticeDFT.Continuum.Box.deriv_eq_zero_of_accPt, LatticeDFT.Continuum.Box.field,
    LatticeDFT.Continuum.Box.lipschitz_field, LatticeDFT.Continuum.Box.eq_zero_of_accPt,
    LatticeDFT.Continuum.Box.ucp_of_isGS}
  \leanok
  \uses{lem:box_ode}
  A ground state vanishes only on a null set.
\end{theorem}
\begin{proof}
  \leanok
  If the zeros accumulated at $x_0\in(0,1)$, then $\psi(x_0)=\psi'(x_0)=0$, and uniqueness for the linear ODE
  $(\psi,\psi')'=(\psi',(v-E)\psi)$ gives $\psi\equiv0$ on $(0,1)$, contradicting $\langle\psi,\psi\rangle>0$. So the
  zeros in each $[1/(k+2),1-1/(k+2)]$ are finite, and all zeros in $(0,1)$ are countable.
\end{proof}

\begin{theorem}[Hohenberg--Kohn for one particle in a box]
  \label{thm:hk_box}
  \lean{LatticeDFT.Continuum.Box.hk_box}
  \leanok
  \uses{thm:hk_continuum, lem:box_setting, thm:box_ucp}
  If ground states of continuous potentials $v$ and $v'$ have the same density $\psi^2=\psi'^2$, then $v-v'$ is
  constant.
\end{theorem}
\begin{proof}
  \leanok
  All hypotheses of Theorem~\ref{thm:hk_continuum} hold.
\end{proof}
